\slajdid{10-s083}
\begin{frame}[label=10-s083]
\gusttekst\setlength{\abovedisplayskip}{3pt}\setlength{\belowdisplayskip}{3pt}\setlength{\abovedisplayshortskip}{2pt}\setlength{\belowdisplayshortskip}{2pt}\begin{columns}[T,onlytextwidth]\column{.29\textwidth}\centering\input{slike/tikz/s082_ecl_diferencijalni_par.tex}\column{.69\textwidth}
\[V_{BE1}-V_{BE2}\ll0\]
Očigledno je da nejednakost može da zadovolji jedino uslovi
\[V_{BE1}\ll0\qquad V_{BE2}=V_{BE}\]
odnosno situacija da je tranzistor T1 zakočen, a da tranzistor T2 vodi. Otpornik $R_{C2}$ i struja $I_E$ se biraju sa uslovom
\[V_{C2}=V_{CC}-R_{C2}I_{C2}=V_{CC}-R_{C2}I_E>V_R\]
tako da tranzistor T2 kada radi ne ide u zasićenje (sva struja strujnog izvora prolazi kroz tranzistor T2 – \underline{zanemarićemo bazne struje tranzistora T1 i T2}). Napon na kolektoru tranzistora T2 je istovremeno i izlazni napon na izlazu $V_{I2}$
\[V_{O2}=V_{C2}=V_{CC}-R_{C2}I_{E2}\]
Napon na izlazu $V_{I1}$ je
\[V_{O1}=V_{C1}=V_{CC}-R_{C1}I_{C1}=V_{CC}\]
\end{columns}
\end{frame}
